% -*- coding: utf-8 -*-
\documentclass[12pt]{article}
\usepackage[utf8]{inputenc}
\usepackage[T1]{fontenc}
\usepackage[portuguese]{babel}
\usepackage[margin=1in]{geometry}
\usepackage{hyperref}
\usepackage{graphicx}
\usepackage{titlesec}
\usepackage{fancyhdr}
\usepackage{enumitem}
\usepackage{array}
\usepackage{longtable}
\usepackage{amsmath}

% Configuração de espaçamento
\setlength{\parindent}{0.5in}
\setlength{\parskip}{0.1in}

% Estilos de títulos
\titleformat{\section}{\Large\bfseries}{\thesection}{0.5em}{}[\vspace{0.2em}]
\titleformat{\subsection}{\large\bfseries}{\thesubsection}{0.5em}{}[\vspace{0.1em}]
\titleformat{\subsubsection}{\normalsize\bfseries}{\thesubsubsection}{0.3em}{}

% Cabeçalho e rodapé
\pagestyle{fancy}
\fancyhf{}
\rhead{SPE UCAN Student Chapter}
\lhead{Documento de Apresentação Oficial}
\cfoot{\thepage}

% Hiperlinks
\hypersetup{
    colorlinks=true,
    linkcolor=blue,
    urlcolor=blue,
    pdftitle={SPE UCAN - Apresentação Oficial},
    pdfauthor={Board Executivo SPE UCAN},
    pdfsubject={Estrutura Organizacional e Programa de Actividades}
}

\title{\textbf{APRESENTAÇÃO OFICIAL DO STUDENT CHAPTER}\\[0.5em]
\large Estrutura Organizacional, Objectivos e Programa de Actividades\\[0.2em]
\normalsize Society of Petroleum Engineers - UCAN}

\author{Board Executivo da SPE UCAN Student Chapter\\
Instituto de Recursos Minerais, Ambiente e Tecnologias (IRMAT)\\
Universidade Católica de Angola (UCAN)}

\date{Luanda, Outubro de 2025}

\begin{document}

\maketitle

\begin{center}
\textit{Este documento apresenta a estrutura, objectivos e programa de actividades\\
da SPE UCAN Student Chapter, oficialmente reconhecida em 15 de Outubro de 2025.}
\end{center}

\newpage
\tableofcontents
\newpage

\section{Enquadramento Institucional e Reconhecimento}

\subsection{Reconhecimento Oficial}

O documento de Estrutura Organizacional, Objectivos e Programa de Actividades da SPE UCAN foi elaborado pelo Board Executivo da SPE UCAN Student Chapter, com supervisão académica do Faculty Advisor do Instituto de Recursos Minerais, Ambiente e Tecnologias (IRMAT), e apoio da SPE Angola Section.

A sua redacção e estrutura seguem as normas estabelecidas pela Society of Petroleum Engineers (SPE International), em conformidade com as políticas internas da Universidade Católica de Angola (UCAN).

O capítulo foi oficialmente reconhecido em 15 de Outubro de 2025, representando um marco significativo para a universidade e para a comunidade académica angolana.

\subsection{Estrutura de Coordenação}

\noindent\textbf{Coordenação Institucional:}
\begin{itemize}
    \item Instituto de Recursos Minerais, Ambiente e Tecnologias (IRMAT)
    \item SPE Angola Section -- Coordenação dos Student Chapters
    \item Reitoria da Universidade Católica de Angola (UCAN)
\end{itemize}

\noindent\textbf{Supervisão e Orientação Académica:}
\begin{itemize}
    \item Director do IRMAT -- Eng.º Tommaso Depippo
    \item Director-adjunto do IRMAT -- Prof. Osvaldo Maiato
    \item Faculty Advisor -- Prof. Emanuel Leopoldo
\end{itemize}

\subsection{Composição do Board Executivo}

\begin{center}
\begin{longtable}{|p{3cm}|p{4cm}|p{4cm}|}
\hline
\textbf{Função} & \textbf{Nome} & \textbf{Curso} \\
\hline
Presidente & Vanessa Cabenguela & Ciências da Terra e Geo-recursos \\
Vice-presidente & Fátima Sahungo & Ciências da Terra e Geo-recursos \\
Secretário & Rosandro Kamahonjo & Engenharia de Petróleo \\
Tesoureira & Prijana Janota & Ciências da Terra e Geo-recursos \\
Membership Chairperson & António Kialongo & Engenharia de Petróleo \\
\hline
\caption{Composição do Board Executivo 2025-2026}
\end{longtable}
\end{center}

\noindent\textbf{1ª Edição -- Outubro de 2025}\\
\textit{Documento interno de carácter académico e institucional. Todos os direitos reservados à SPE UCAN Student Chapter.}\\
\textit{Luanda, aos 31 de Outubro de 2025}

\newpage

\section{Contexto da Indústria Petrolífera Angolana}

Angola é o segundo maior produtor de petróleo em África Subsaariana, com produção de 1,4 milhões de barris por dia em 2025. O sector representa 90\% das exportações e 30\% do PIB, mas enfrenta desafios como dependência económica, flutuações de preços e transição para energias renováveis.

\subsection{História da Indústria Petrolífera em Angola}

\begin{itemize}
    \item \textbf{Década de 1950:} Descobertas iniciais de petróleo no território angolano
    \item \textbf{Pós-independência (1975):} Nacionalização com criação da Sonangol
    \item \textbf{Boom offshore (1990s-2000s):} Descoberta de campos como Girassol e Dalia
    \item \textbf{Actualidade:} Foco em deepwater e desenvolvimento de gás natural
    \item \textbf{Transição energética (2020-2025):} Investimentos em energias renováveis
\end{itemize}

\subsection{Principais Operadores em Angola}

\begin{itemize}
    \item \textbf{Sonangol:} Empresa estatal, controla 40\% da produção, opera blocos 0, 14, 17
    \item \textbf{Chevron:} Responsável pelo bloco BBLT, produz cerca de 200.000 bpd
    \item \textbf{TotalEnergies:} Operadora dos blocos Kaombo e Norte
    \item \textbf{BP e ExxonMobil:} Envolvidas em LNG e bloco 31
    \item \textbf{Investidores chineses:} Sinopec e CNOOC com presença significativa
\end{itemize}

\subsection{Desenvolvimentos Recentes (2020-2025)}

\begin{itemize}
    \item Produção aumentou de 1,1 milhões bpd (2020) para 1,4 milhões (2025)
    \item Novos leilões atraíram investimentos de 10 bilhões USD
    \item Projecto Angola LNG com capacidade de 5 milhões de toneladas
    \item Desenvolvimento de parques solares em Benguela e Namibe
    \item Ênfase em sustentabilidade e compliance ambiental
\end{itemize}

\subsection{Reservas e Recursos}

\begin{itemize}
    \item Reservas provadas: 9,1 bilhões de barris
    \item Gás natural: 330 bilhões de metros cúbicos
    \item \textbf{Campos principais:}
    \begin{itemize}
        \item Girassol: 800 milhões barris
        \item Dalia: 600 milhões barris
        \item Pazflor: 500 milhões barris
    \end{itemize}
\end{itemize}

\newpage

\section{Missão, Visão e Objectivos da SPE UCAN}

\subsection{Declaração de Missão}

A missão da SPE UCAN consiste em promover a excelência técnica e científica, incentivar a ética e a liderança, e criar oportunidades de desenvolvimento pessoal e profissional para os seus membros. O capítulo opera como ponte entre a academia e a indústria petrolífera nacional e internacional.

\subsection{Declaração de Visão}

A visão do capítulo é tornar-se uma referência nacional e regional em qualidade académica, inovação e ligação com a indústria energética, consolidando o papel da UCAN e do IRMAT como centros de excelência e produção de conhecimento.

\subsection{Objectivos Estratégicos}

O capítulo nasce orientado pelos seguintes objectivos:

\begin{enumerate}
    \item \textbf{Fortalecer a formação prática} dos estudantes nas disciplinas de engenharia e geociências
    \item \textbf{Fomentar a investigação aplicada} em colaboração com a indústria
    \item \textbf{Estabelecer parcerias} com empresas e instituições do sector energético
    \item \textbf{Garantir uma gestão participativa, transparente e sustentável} do capítulo
    \item \textbf{Promover a diversidade e inclusão} entre os membros
    \item \textbf{Facilitar intercâmbios académicos} internacionais
    \item \textbf{Contribuir para o desenvolvimento} sustentável de Angola
\end{enumerate}

\subsection{Princípios de Gestão}

O SPE UCAN opera segundo os seguintes princípios:

\begin{itemize}
    \item Excelência técnica e científica
    \item Ética profissional e integridade
    \item Transparência nas decisões
    \item Participação democrática
    \item Responsabilidade social e ambiental
    \item Inovação contínua
    \item Inclusão e igualdade
\end{itemize}

\newpage

\section{Oportunidades Profissionais para Estudantes}

Estudantes da UCAN têm acesso único à indústria petrolífera angolana, com oportunidades significativas em termos de carreira, desenvolvimento técnico e remuneração competitiva.

\subsection{Perspectivas de Carreira}

Os salários médios para recém-graduados situam-se entre 60.000 a 100.000 Kz mensais, enquanto profissionais com experiência podem auferir 150.000 Kz ou mais, dependendo da função e da empresa.

\subsection{Áreas de Especialização}

\begin{enumerate}
    \item \textbf{Engenharia de Reservatórios:} Análise, modelagem e otimização de campos petrolíferos
    \begin{itemize}
        \item Salário inicial: 60.000 Kz
        \item Empresas: Sonangol P\&P, Chevron
        \item Competências: Simulação, análise de dados, modelagem numérica
    \end{itemize}

    \item \textbf{Engenharia de Produção:} Supervisão de extracção, manutenção de plataformas e optimização operacional
    \begin{itemize}
        \item Salário inicial: 70.000 Kz
        \item Oportunidades em operações offshore deepwater
        \item Competências: Gestão de projectos, segurança operacional
    \end{itemize}

    \item \textbf{Engenharia Ambiental e HSE:} Foco em sustentabilidade, redução de emissões e compliance regulatório
    \begin{itemize}
        \item Crescente demanda devido a regulamentações ESG
        \item Competências: Análise ambiental, gestão de riscos
    \end{itemize}

    \item \textbf{Engenharia de Refino e Processamento:} Processamento de petróleo em refinarias
    \begin{itemize}
        \item Localização: Luanda Refinery, futuras instalações
    \end{itemize}

    \item \textbf{Geociências e Exploração:} Prospecção, avaliação de recursos e interpretação geológica
    \begin{itemize}
        \item Carreira: Explorador, geólogo estrutural, geofísico
    \end{itemize}
\end{enumerate}

\subsection{Oportunidades de Estágio}

\begin{itemize}
    \item \textbf{Sonangol:} Programas de estágio de 12 meses para recém-graduados
    \item \textbf{Chevron Angola:} Internships de 6 a 12 meses
    \item \textbf{TotalEnergies:} Graduate programmes com duração de 24 meses
    \item \textbf{Outras operadoras:} BP, ExxonMobil, Azule Energy
\end{itemize}

\subsection{Programas de Pesquisa e Desenvolvimento}

\begin{itemize}
    \item Captura e armazenamento de carbono (CCUS)
    \item Inteligência artificial e machine learning em operações
    \item Energias renováveis e transição energética
    \item Eficiência operacional e otimização de custos
    \item Segurança offshore e resiliência operacional
\end{itemize}

\newpage

\section{Mensagem de Agradecimentos}

A criação da SPE UCAN Student Chapter é fruto de um esforço conjunto que uniu entusiasmo estudantil, visão institucional e orientação profissional. O seu estabelecimento oficial representa não apenas um marco académico para o Instituto de Recursos Minerais, Ambiente e Tecnologias (IRMAT) e para a Universidade Católica de Angola (UCAN), mas também o resultado da dedicação de pessoas que acreditaram na importância de inserir a universidade na comunidade internacional da Society of Petroleum Engineers (SPE).

Expressamos o nosso profundo agradecimento por todo o apoio académico, institucional e logístico que tornou possível a criação deste capítulo estudantil.

Estendemos igualmente os nossos agradecimentos à SPE Angola Section pela orientação, acompanhamento e validação de todo o processo de estruturação. Dirigimos um especial agradecimento ao Eng.º Abel Cruz, Director da SPE Angola Section, e o nosso profundo apreço estende-se ao Eng.º Délcio Falcão, Coordenador dos Núcleos Universitários da SPE Angola Section.

Um agradecimento especial à Eng.ª Teresa Júlio, cuja iniciativa foi essencial no contacto inicial com a SPE International e com o Faculty Advisor, Eng.º Emanuel Leopoldo, ao qual manifestamos também a nossa sincera gratidão.

Aos 25 estudantes fundadores, o nosso sincero reconhecimento pelo empenho e disciplina demonstrado durante todo o percurso de estabelecimento do SPE UCAN Student Chapter.

\begin{quote}
\textit{``Começamos engatinhando. Chegou a hora de levantar, mantermo-nos em pé e aprendermos a andar e depois correr. Tornemos a SPE UCAN a vanguarda académica que abre caminho para os outros, inova na prática e lidera pela excelência, identidade única da Mater et Magistra - UCAN''}
\end{quote}

\newpage

\section{Conclusão}

Este relatório institucional estabelece os fundamentos académicos, administrativos e estratégicos da SPE UCAN Student Chapter. Mais do que um documento formal, representa a declaração de um compromisso profundo: o de formar uma nova geração de engenheiros, cientistas e líderes éticos preparados para enfrentar os desafios do desenvolvimento energético de Angola e do mundo.

A SPE UCAN Student Chapter nasce como um espaço de convergência entre o saber, a inovação e a juventude universitária, promovendo a energia do conhecimento, o rigor técnico e o espírito de colaboração que caracterizam a Universidade Católica de Angola e a Society of Petroleum Engineers.

Com dedicação à excelência, esperamos consolidar-nos como referência nacional e internacional no panorama académico e profissional do sector energético.

\begin{flushright}
\textit{Luanda, Outubro de 2025}\\[0.5cm]
\textbf{Board Executivo SPE UCAN Student Chapter}
\end{flushright}

\newpage

\section*{Tabela de Siglas e Abreviaturas}

\begin{center}
\begin{longtable}{|p{3cm}|p{9cm}|}
\hline
\textbf{Sigla} & \textbf{Descrição} \\
\hline
AIME & American Institute of Mining, Metallurgical, and Petroleum Engineers \\
\hline
bpd & Barris por dia (barrels per day) \\
\hline
CCUS & Captura, Utilização e Armazenamento de Carbono \\
\hline
IRMAT & Instituto de Recursos Minerais, Ambiente e Tecnologias \\
\hline
HSE & Health, Safety and Environment (Segurança, Saúde e Ambiente) \\
\hline
LNG & Liquefied Natural Gas (Gás Natural Liquefeito) \\
\hline
MIREMPET & Ministério dos Recursos Minerais, Petróleo e Gás \\
\hline
OPEC & Organization of the Petroleum Exporting Countries \\
\hline
SPE & Society of Petroleum Engineers \\
\hline
SPE UCAN & Society of Petroleum Engineers - UCAN Student Chapter \\
\hline
UCAN & Universidade Católica de Angola \\
\hline
\end{longtable}
\end{center}

\end{document}
